\section{Ambient one-point extensions}\label{sec:ambient}

\subsection{The hereditary theorem}
Let \(\lambda:R\to S=R_\Sigma\) be a stably flat universal localisation of finite-dimensional algebras at a finite set \(\Sigma\) of maps between finitely generated projective left modules, with \(T=\bigoplus_\sigma\Cone(\sigma)\), \(B\), \(F\), \(\Theta\) as in \Cref{sec:preliminaries}, and bounded profile \(\Prof(\lambda)=\{Y\in D(B):FY\in\Db(R\text{-}\modcat)\}\subseteq\Dfd(B)\).  Let \(W\) be a finite-dimensional right \(R\)-module and form the one-point extension
\[
 A'=R[W]=\begin{pmatrix}k&W\\0&R\end{pmatrix},\qquad p=e_{11},\qquad e=e_{22}.
\]
Left \(A'\)-modules are triples \((V,Z,\varphi)\) with \(\varphi:W\otimes_RZ\to V\); the corner functor \(j^*=e(-)\) has derived left adjoint \(j_!Z=A'e\Ltensor_RZ=(W\Ltensor_RZ,Z,\id)\), which is fully faithful, coproduct-preserving and compact-preserving, and \(i_*:D(k)\to D(A')\) is restriction along \(A'\to A'/A'eA'=k\).  Put
\[
 \Sigma'=\{A'e\otimes_R\sigma:\sigma\in\Sigma\}=j_!\Sigma .
\]

\begin{theorem}[Ambient one-point extension]\label{thm:ambient}
With the notation above:
\begin{enumerate}[label=\textup{(\alph*)}]
\item\label{amb:a} The universal localisation of \(A'\) at \(\Sigma'\) is
\begin{equation}\label{eq:ambient-target}
 S'=A'_{\Sigma'}\cong\begin{pmatrix}k&W\otimes_RS\\0&S\end{pmatrix},
\end{equation}
with \(\lambda':A'\to S'\) given by \(\lambda\) on \(R\), \(w\mapsto w\otimes1\) on \(W\), and the identity on \(k\).  It is stably flat if and only if
\begin{equation}\label{eq:tor-condition}
 \Tor_i^R(W,{}_RS)=0\qquad(i\ge1),
\end{equation}
which holds for every \(W\) when \({}_RS\) is projective.
\item\label{amb:b} Under \eqref{eq:tor-condition}, \(\cK_{\lambda'}=j_!(\cK_\lambda)\), \(\Kc_{\lambda'}=\thick(j_!T)\), and \(\RHom_R(T,T)\to\RHom_{A'}(j_!T,j_!T)\) is a quasi-isomorphism of dg algebras; the kernel dg algebra of \(\lambda'\), marked by \(j_!T\), is \(B\).
\item\label{amb:c} Under \eqref{eq:tor-condition}, in the marking of \ref{amb:b},
\begin{equation}\label{eq:profile-intersection}
 \Prof(\lambda')=\Prof(\lambda)\cap\cU_{\Theta(D_kW)}=\{Y\in\Prof(\lambda):W\Ltensor_RFY\in\Db(k\text{-}\modcat)\},
\end{equation}
so that \(\DeltaB(\lambda')\simeq\bigl(\Prof(\lambda)\cap\cU_{\Theta(D_kW)}\bigr)/\per B\).  At the level of detectors, \(D_k(j_!T)\simeq D_kT\oplus\Theta(D_kW)\) in \(D(B)\).
\item\label{amb:d} If \(\Prof(\lambda)=\Dfd(B)\) and \({}_RS\) is projective, then for every \(E\in\Dfd(B)\) there is a finite right \(R\)-module \(W\) with \(\Prof(\lambda')=\cU_E\).
\end{enumerate}
\end{theorem}

\begin{proof}
\ref{amb:a}, \emph{universal property.}  The displayed map inverts \(\Sigma'\), since extension along it restricted to the corner is extension along \(\lambda\).  Let \(f:A'\to D\) be a ring map inverting \(\Sigma'\), and put \(p_D=f(p)\), \(e_D=f(e)\).  Applying the exact functor \(e_D(-)\) to the inverted maps shows that the corner map \(R\to e_DDe_D\) inverts \(\Sigma\), so it extends uniquely to \(g:S\to e_DDe_D\).  The restriction of \(f\) to \(W\subseteq p_DDe_D\) extends to \(W\otimes_RS\to p_DDe_D\), \(w\otimes s\mapsto f(w)g(s)\), which is well defined because \(f(wr)=f(w)f(r)\) and \(g\) extends \(f|_R\).  Together with \(f\) on the corner \(k\), these define a ring map \(S'\to D\) extending \(f\); balancedness is the only nontrivial multiplicativity check, and uniqueness holds because \(k\), \(W\) and \(R\) generate \(A'\).  This proves \eqref{eq:ambient-target} independently of \eqref{eq:tor-condition}.

\ref{amb:a}, \emph{stable flatness.}  Let \(L=S\Ltensor_R-\) be the Bousfield localisation of \(D(R)\) with kernel \(\cK_\lambda\), so that \(LR=S\) as an \(R\)-module.  Since \(j_!\) preserves coproducts, cones and summands and is fully faithful, the localising subcategory generated by the new cones is
\begin{equation}\label{eq:new-kernel}
 \cC'=\Loc(\Cone\Sigma')=j_!(\cK_\lambda),
\end{equation}
which is generated by compact objects; let \(L'\) be the associated smashing localisation of \(D(A')\).  For \(Z\in D(R)\) apply \(j_!\) to the localisation triangle \(K_Z\to Z\to LZ\): the first term lies in \(\cC'\), and \(j_!LZ\) is \(\cC'\)-local because \(\Hom_{D(A')}(j_!K,j_!LZ[n])\cong\Hom_{D(R)}(K,LZ[n])=0\) for \(K\in\cK_\lambda\), by \(j^*j_!\cong\id\).  Hence
\begin{equation}\label{eq:L-prime-j}
 L'(j_!Z)\simeq j_!(LZ).
\end{equation}
Similarly \(i_*k=(k,0,0)\) is \(\cC'\)-local, since \(\Hom(j_!K,i_*k[n])\cong\Hom(K,j^*i_*k[n])=0\), so \(L'(i_*k)=i_*k\).  As \(A'=A'p\oplus A'e=i_*k\oplus j_!R\),
\begin{equation}\label{eq:L-prime-A}
 L'(A')\simeq i_*k\oplus j_!(S)=\bigl(k\oplus(W\Ltensor_RS),\ S,\ \ldots\bigr),
\end{equation}
whose cohomology is \(H^0(L'A')={}_{A'}S'\) and \(H^{-i}(L'A')=i_*\Tor_i^R(W,S)\) for \(i\ge1\), with \(H^{>0}=0\).

A smashing localisation of the derived category of an algebra \(A\) is induced by a homological ring epimorphism \(A\to A''\) with \(A''\) an ordinary algebra if and only if the localised regular object \(LA\) has cohomology concentrated in degree zero; in that case \(A''=\End_{D(A)}(LA)^{\op}\) in the left-module convention.  Indeed, Nicolás--Saorín \cite[Theorem~5]{NicolasSaorin09} model the smashing localisation by a homological epimorphism to a dg algebra \(\mathcal C\) whose restricted regular module is \(LA\); if \(H^{i}(LA)=0\) for \(i\ne0\), the maps of dg algebras \(\mathcal C\leftarrow\tau_{\le0}\mathcal C\to H^0(\mathcal C)\) under \(A\) are quasi-isomorphisms, replacing the dg target by an ordinary algebra with the same fully faithful restriction; conversely the restriction of the regular module of an ordinary homological target is a stalk complex.  The same reduction appears in the commutative setting in \cite[proof of Theorem~4.9]{AMSTV20}.

Under \eqref{eq:tor-condition}, \eqref{eq:L-prime-A} is concentrated in degree \(0\).  Write \(I=i_*k\) and \(J=j_!S\).  Then \(\End(I)=k\), \(\End(J)=\End_{D(R)}(S)=S^{\op}\) by homologicality of \(\lambda\), \(\Hom(J,I)=\Hom_{D(R)}(S,j^*i_*k)=0\), and \(\Hom(I,J)=pJ=W\otimes_RS\) because \(I=A'p\) is projective.  Hence \(\End_{D(A')}(L'A')^{\op}\cong S'\), and applying \(L'\) to right multiplications on \(A'\) recovers the map \(\lambda'\) of \eqref{eq:ambient-target}.  Thus \(\lambda'\) is a homological ring epimorphism, and being a universal localisation it is stably flat.  Conversely, if \(\lambda'\) were stably flat while some \(\Tor_i^R(W,S)\ne0\), Neeman--Ranicki \cite{NeemanRanicki04} would identify its kernel with \eqref{eq:new-kernel} and its localised regular object with the stalk module \(S'\), contradicting \eqref{eq:L-prime-A}.

\ref{amb:b}.  By stable flatness and \Cref{lem:compact-intersection}, \(\cK_{\lambda'}=\cC'\) and \(\Kc_{\lambda'}=\thick\{j_!\Cone\sigma\}=\thick(j_!T)\).  On K-projective representatives, derived adjunction and \(j^*j_!\cong\id\) give a quasi-isomorphism of dg Hom complexes \(\RHom_R(T,T)\to\RHom_{A'}(j_!T,j_!T)\) compatible with composition; taking graded opposites identifies the kernel dg algebras.

\ref{amb:c}.  A complex of \(A'\)-modules has bounded finite-dimensional cohomology iff both corner complexes do, since the two exact idempotent functors \(p(-)\), \(e(-)\) have underlying vector spaces summing to that of the complex.  By the formula for \(j_!\), \(j_!FY\) is bounded iff \(FY\) and \(W\Ltensor_RFY\) are, giving the second description in \eqref{eq:profile-intersection}.  Tensor duality and adjunction give, with all sides fixed as in \Cref{sec:preliminaries},
\[
 D_k\bigl(W\Ltensor_RFY\bigr)\simeq\RHom_R(FY,D_kW)\simeq\RHom_B\bigl(Y,\Theta(D_kW)\bigr),
\]
and \(\Theta(D_kW)\) is proper because \(T\) is a bounded complex of finitely generated projectives.  This proves \eqref{eq:profile-intersection}; the quotient formula is \Cref{cor:relative-dg-bounded} for \(\lambda'\), and the detector formula follows from the two components of \(j_!T\) and \eqref{eq:tensor-duality}.

\ref{amb:d}.  Let \(E\in\Dfd(B)\) and \(X=FE\), which lies in \(\Db(R\text{-}\modcat)\) by fullness.  \Cref{lem:module-detector} gives a module \(N\) with \(\RHom_R(FY,X)\) bounded iff \(\RHom_R(FY,N)\) bounded, for all \(Y\in\Dfd(B)\).  Put \(W=D_kN\), so \(D_kW=N\); condition \eqref{eq:tor-condition} holds because \({}_RS\) is projective.  Full faithfulness of \(F\) gives \(\RHom_B(Y,E)\simeq\RHom_R(FY,X)\), hence \(\cU_E=\cU_{\Theta(D_kW)}\), and \ref{amb:c} with \(\Prof(\lambda)=\Dfd(B)\) gives \(\Prof(\lambda')=\cU_E\).
\end{proof}

\begin{remark}
Formula \eqref{eq:profile-intersection} is an intersection: ambient one-point extensions can only shrink a profile, and \(W=0\) preserves it.  The projectivity hypothesis in \ref{amb:d} is on \({}_RS\), the side paired with \(W\) in \eqref{eq:tor-condition}; projectivity of \(S_R\) would not suffice.  When \eqref{eq:tor-condition} fails, \(\lambda'\) is a universal localisation of finite-dimensional algebras that is not stably flat, a source of such examples.
\end{remark}

\subsection{Marks--Nakayama kernels}
Let \(k\) be algebraically closed, \(d\ge3\), \(A_d=k\widetilde Q_{2d}/\mathfrak r^d\), and let \(J\) be Marks-admissible as in (\Cref{prop:marks-projectivity,thm:nakayama-intrinsic}).  Following \Cref{sec:nakayama} we work with right \(A_d\)-modules, that is, with \(R=A_d^{\op}\); the kernel dg algebra is \(B_J\), the profile is full, and \(S_J\) is projective on both sides.

\begin{corollary}[Profiles with Nakayama kernel]\label{cor:nakayama-profiles}
For every Marks-admissible \(J\), the marked bounded defects realised by stably flat universal localisations of finite-dimensional algebras with kernel dg algebra \(B_J\) are exactly
\[
 \{\cU_E/\per B_J:E\in\Dfd(B_J)\}.
\]
\end{corollary}

\begin{proof}
Necessity is \Cref{cor:pair-invariant}.  Sufficiency is \Cref{thm:ambient}\ref{amb:d}, whose hypotheses are supplied by (\Cref{prop:marks-projectivity,thm:nakayama-intrinsic}).
\end{proof}

We compute these defects explicitly in small rank.  Write \(M(i,l)=e_iA_d/\rad^l(e_iA_d)\), \(1\le l<d\), for the uniserial right module with composition factors \(S_i,\dots,S_{i+l-1}\), indices modulo \(2d\).

\begin{lemma}[Stable Hom in \(\underline{\modcat}\,A_d\)]\label{lem:stable-hom}
\begin{enumerate}[label=\textup{(\roman*)}]
\item \(\Omega M(i,l)=M(i+l,d-l)\), \(M(i,l)[1]=M(i+l-d,d-l)\), and \(\Omega^4M(i,l)\cong M(i,l)\).
\item For \(X=M(i,l)\), \(Y=M(j,m)\) and \(r\equiv i-j\pmod{2d}\) with \(0\le r\le 2d-1\),
\[
 \dim_k\underline{\Hom}(X,Y)=\begin{cases}1,&r<m\le r+l<d,\\0,&\text{otherwise.}\end{cases}
\]
\item For \(X,Y\in\Db(A_d^{\op}\text{-}\modcat)\), \(\Hom_{D}(X,Y[n])\cong\Hom_{\Dsg}(qX,qY[n])\) for \(n\gg0\).  Hence, for a finite module \(N\),
\begin{equation}\label{eq:profile-stable}
 \Ext^{\gg0}(X,N)=0\iff\underline{\Hom}(qX,qN[t])=0\ \text{for all }t\in\bbZ .
\end{equation}
\end{enumerate}
\end{lemma}

\begin{proof}
(i) The kernel of the projective cover of \(M(i,l)\) is the uniserial module starting at \(i+l\) of length \(d-l\); solving for the inverse syzygy gives the shift formula; two syzygies add \(d\) to the top index and four add \(2d\).

(ii) A map \(X\to Y\) is determined by the image of the top of \(X\), which can only lie in the layer of \(Y\) indexed \(r\); such a layer exists iff \(r<m\), and the image is annihilated by paths of length \(l\) iff \(m\le r+l\).  The map factors through the projective cover \(P_j\to Y\) exactly when its image lifts to layer \(r\) of \(P_j\) with the same annihilation relation, that is, iff \(r+l\ge d\); every map through a projective factors through \(P_j\).  Removing these maps gives the formula, over every field.

(iii) Replace \(X\) and \(Y\), modulo bounded projective complexes, by shifted finite modules; since \(A_d\) is self-injective these complexes are also bounded injective, and Hom into or out of them vanishes in large positive degree.  For modules, \(\Ext^n(M,N)\cong\underline\Hom(\Omega^nM,N)=\Hom_{\Dsg}(qM,qN[n])\) for \(n\ge1\).  Periodicity (i) gives \eqref{eq:profile-stable}.
\end{proof}

\begin{proposition}[Small lattices]\label{prop:m3}
Let \(\cQ_J=\Dfd(B_J)/\per B_J\simeq\thick_{\underline{\modcat}A_d}\{S_i:i\in J\cup(J+d)\}\) \textup{(\Cref{prop:marks-projectivity,thm:nakayama-intrinsic})}.
\begin{enumerate}[label=\textup{(\roman*)}]
\item\label{m3:single} If \(J=\{a\}\) is a singleton (all admissible \(J\) for \(d=3\)), \(\cQ_J\) is the split four-periodic category on the shifts of \(S_a\), its thick subcategories are \(0\) and \(\cQ_J\), and both are realised: by \(W=0\) and by \(W=D_kS_a\).
\item\label{m3:opposite} If \(d=4\) and \(J=\{a,a+2\}\), then \(\cQ_J=O_a\oplus O_{a+2}\) with \(O_i\) the split four-periodic category on the shifts of \(S_i\); its thick subcategories form the four-element Boolean lattice and all are realised, the subcategory omitting \(O_i\) by \(W=D_kS_i\) and zero by \(W=D_k(S_a\oplus S_{a+2})\).
\item\label{m3:adjacent} If \(d=4\) and \(J=\{a,a+1\}\), put \(A_0=S_a\), \(A_1=S_{a+1}\), \(A_2=M(a,2)\) and \(O_r=\add\{A_r[t]:t\in\bbZ\}\).  The uniserial extension gives a non-split triangle
\begin{equation}\label{eq:m3-triangle}
 A_1\longrightarrow A_2\longrightarrow A_0\longrightarrow A_1[1],
\end{equation}
\(\cQ_J=\add(O_0,O_1,O_2)\), and the thick subcategories of \(\cQ_J\) are exactly
\[
 0,\quad O_0,\quad O_1,\quad O_2,\quad\cQ_J,
\]
forming the lattice \(M_3\).  All five are realised: \(\cQ_J\) by \(W=0\), \(O_0\) by \(W=D_kM(a,2)\), \(O_1\) by \(W=D_kS_a\), \(O_2\) by \(W=D_kS_{a+1}\), and \(0\) by \(W=D_k(S_a\oplus S_{a+1})\).
\end{enumerate}
In each case the poset of realisable bounded defects with kernel \(B_J\) is the full lattice of thick subcategories of \(\cQ_J\).
\end{proposition}

\begin{proof}
By \Cref{cor:nakayama-profiles} and \eqref{eq:profile-stable}, the realisable defects are the images of the subcategories \(\{X:\underline\Hom(qX,qN[t])=0\ \forall t\}\) for finite modules \(N\), and every detector \(E\) may be replaced by such an \(N\) through \Cref{lem:module-detector}.  In all cases every indecomposable has endomorphism ring \(k\) in the stable category and each orbit has four objects, by \Cref{lem:stable-hom}(i),(ii).

\ref{m3:single}: all off-diagonal stable Homs between the four shifts of \(S_a\) vanish by \Cref{lem:stable-hom}(ii) (the case \(d=3\), \(a=1\) is \Cref{lem:a63-stable}); every morphism between finite sums reduces to identities and zeros, cones stay in the additive category, and a nonzero thick subcategory contains an indecomposable and hence, by shift closure, all four.  \(W=0\) realises \(\cQ_J\); \(N=S_a\) excludes every indecomposable since \(\underline\Hom(S_a[t],S_a[t])\ne0\), realising \(0\).

\ref{m3:opposite}: \Cref{lem:stable-hom}(ii) gives \(\underline\Hom(O_a,O_{a+2}[t])=0=\underline\Hom(O_{a+2},O_a[t])\) for all \(t\), so \(\cQ_J\) is the orthogonal sum of two split categories; thick subcategories are sums of a subset of the two orbits.  For \(N\) the sum of the simples of the omitted orbits, \eqref{eq:profile-stable} picks out exactly the remaining orbits.

\ref{m3:adjacent}: \Cref{lem:stable-hom}(ii) gives the stable Hom table below, rows indexed by the source, columns by the target, entries listing dimensions for target shifts \(0,1,2,3\):
\[
\begin{array}{c|ccc}
 & A_0 & A_1 & A_2\\ \hline
 A_0 & (1,0,0,0) & (0,1,0,0) & (0,0,0,0)\\
 A_1 & (0,0,0,0) & (1,0,0,0) & (1,0,0,0)\\
 A_2 & (1,0,0,0) & (0,0,0,0) & (1,0,0,0)
\end{array}
\]
Thus among the twelve indecomposables the only non-isomorphisms are scalar multiples of the arrows
\(A_1[t]\to A_2[t]\to A_0[t]\to A_1[t+1]\), the rotations of \eqref{eq:m3-triangle}, and consecutive compositions vanish because they are consecutive maps in a triangle; the radical of the additive category \(\add(O_0,O_1,O_2)\) has square zero.  Given a morphism between finite sums, split off identities by block row and column operations until all entries are radical; then for each source type the only nonzero block maps onto its unique successor, and independent base changes reduce it to rank normal form.  Every morphism is therefore a direct sum of identities, zeros, and copies of the arrows, whose cones are the third terms of rotations of \eqref{eq:m3-triangle}.  Hence \(\add(O_0,O_1,O_2)\) is thick, contains \(S_a,S_{a+1}\), and equals \(\cQ_J\).  A nonzero thick subcategory contains a complete orbit; any two orbits force the third by \eqref{eq:m3-triangle}; each single orbit is split and thick by the table.  This is the lattice \(M_3\).  The realising modules follow from the columns of the table via \eqref{eq:profile-stable}: \(N=S_a\) excludes \(O_0,O_2\) and keeps \(O_1\); \(N=S_{a+1}\) keeps \(O_2\); \(N=M(a,2)\) keeps \(O_0\); \(N=S_a\oplus S_{a+1}\) excludes everything.  The stable Hom tables for \(d=3,4\) are checked exactly by \texttt{verify\_nakayama\_profiles.py} (\Cref{app:computation}).
\end{proof}

\begin{corollary}[No support space for adjacent pairs]\label{cor:non-distributive}
For \(d=4\) and \(J=\{a,a+1\}\) the poset of realisable bounded defects with kernel \(B_J\) is the non-distributive lattice \(M_3\).  It is therefore neither order-isomorphic nor order-anti-isomorphic to the lattice of closed subsets of any topological space.  In particular no analogue of the support classification of \Cref{thm:support-classification} exists for this kernel.
\end{corollary}

\begin{proof}
Closed subsets of a space form a distributive lattice under union and intersection; distributivity is preserved by order isomorphisms and anti-isomorphisms; \(M_3\) is not distributive.
\end{proof}

\subsection{A genuinely dg boundary}

\begin{theorem}[The twenty-dimensional boundary]\label{thm:dg-boundary}
Let \(k\) be algebraically closed, \(A=A_3=k\widetilde Q_6/\mathfrak r^3\) with right modules as in \Cref{subsec:a63}, \(R=A^{\op}\), and \(\lambda\) the Marks localisation at \(\{S_1,S_4\}\) with kernel dg algebra
\[
 B^0=ke_1\oplus ke_4,\qquad B^{-1}=ka\oplus kb,\qquad a=e_4ae_1,\quad b=e_1be_4,\quad ab=ba=0,\quad d=0,
\]
and simple dg modules \(L_1,L_4\) \textup{(\Cref{prop:marks-projectivity,thm:nakayama-intrinsic})}.  Let \(W=D_kS_1\) and \(A'=R[W]\), an algebra of dimension \(1+1+18=20\), and let \(\lambda'\) be the induced localisation.  Then \(\lambda'\) is stably flat with kernel dg algebra \(B\),
\begin{equation}\label{eq:ext-degrees}
 \Ext^q_B(L_1,L_1)=\begin{cases}k,&q\equiv0\pmod4,\ q\ge0,\\0,&\text{otherwise},\end{cases}
 \qquad
 \Ext^q_B(L_4,L_1)=\begin{cases}k,&q\equiv2\pmod4,\ q\ge0,\\0,&\text{otherwise},\end{cases}
\end{equation}
\(\Theta(S_1)=L_1\), \(\cU_{L_1}=\per B\), and
\begin{equation}\label{eq:dg-boundary}
 \DeltaB(\lambda')=0,\qquad \Dfd(B)/\per B\simeq\vect_k^{\bbZ/4}\ne0 .
\end{equation}
Thus the intrinsic connective formula \(\DeltaB\simeq\Dfd(B)/\per B\) fails for a marked kernel dg algebra with nonzero negative cohomology.
\end{theorem}

\begin{proof}
The target of \(\lambda\) is projective on both sides and the profile is full (\Cref{prop:marks-projectivity,thm:nakayama-intrinsic}), so \Cref{thm:ambient} applies with the marked kernel preserved.  The identifications \(F(L_1)=S_1\), \(F(L_4)=S_4\) are \Cref{prop:nakayama-cancellation,prop:a63-bounded-kernel}; full faithfulness gives \(\Theta(S_1)=L_1\), so the new detector is \(\Theta(D_kW)=L_1\).

The short exact sequences \(0\to L_4[1]\to Be_1\to L_1\to0\) and \(0\to L_1[1]\to Be_4\to L_4\to0\) of \Cref{prop:a63-dga} splice to a semi-free resolution of \(L_1\) whose underlying graded module is
\[
 P^{\#}=Be_1\oplus Be_4[2]\oplus Be_1[4]\oplus Be_4[6]\oplus\cdots,
\]
with free generators \(x_n\) of degree \(-2n\) at vertex \(1\) for \(n\) even and at vertex \(4\) for \(n\) odd, and differential \(dx_0=0\), \(dx_n=ax_{n-1}\) for \(n\) odd, \(dx_n=bx_{n-1}\) for \(n\) even.  Each connecting shift is by \(2\), one from the internal degree of the radical generator and one from the cone.  Since \(ab=ba=0\), \(d^2=0\), the cycles are \(x_0\) and the elements \(ax_n,bx_n\), all of which except \(x_0\) are boundaries, so the augmentation to \(L_1\) is a quasi-isomorphism.  The radical acts by zero on \(L_1\), so \(\Hom_B(P,L_1)\) has zero differential, and \(\Hom_B(Be_v[2n],L_1)=e_vL_1\) in degree \(2n\).  Generators at vertex \(1\) occur for \(n=2m\), giving degrees \(4m\); this proves the first formula in \eqref{eq:ext-degrees}.  The resolution of \(L_4\) is obtained by interchanging the vertices and arrows, and its vertex-\(1\) generators occur for \(n=2m+1\), giving degrees \(4m+2\).

For \(\cU_{L_1}=\per B\), let \(Y\in\Dfd(B)\).  Full faithfulness gives \(\RHom_B(Y,L_1)\simeq\RHom_R(FY,S_1)\), and by \Cref{lem:stable-hom}(iii) its cohomology in large degrees is \(\underline\Hom(qFY,qS_1[n])\).  By \Cref{thm:a63-comparison}, \(q\Kb(\lambda)\) is the split four-periodic category on \(qS_1\) with \(\dim\Hom(Y_r,Y_s[n])=1\) iff \(s+n-r\equiv0\pmod4\); every nonzero object has a nonzero map to \(qS_1[n]\) for some \(n\), recurring in all degrees \(n+4m\).  Hence \(\RHom_B(Y,L_1)\) is bounded iff \(qFY=0\) iff \(FY\in\per R\) iff \(Y\in\per B\).  \Cref{thm:ambient}\ref{amb:c} gives \(\Prof(\lambda')=\per B\) and \eqref{eq:dg-boundary}; the intrinsic quotient is the displayed category by \Cref{thm:a63-comparison}.  Both formulas in \eqref{eq:ext-degrees} and the four-periodic Hom pattern are checked exactly by \texttt{verify\_dg\_boundary.py} (\Cref{app:computation}).
\end{proof}

\begin{remark}
The zero defect is reached without the kernel algebra being ordinary: this is the phenomenon of \Cref{thm:dual-numbers} for a kernel dg algebra with \(H^{-1}\ne0\).  The marked dg algebra \(B\) is not quasi-isomorphic to an ordinary algebra; no assertion is made about other objects of its dg Morita class.  As in \Cref{sec:support}, the failure lies inside the finite-dimensional silting range of \Cref{cor:finite-dimensional-precompletion}, so \(\DeltaB(\lambda')=\Defsg(\lambda')=0\) and no summand obstruction is involved.
\end{remark}
